% 【本文件写什么】impl 实现层：ringbuf
%   - crate 路径：os/components/wateros-klog/klog-impl/impl-ringbuf/src/
% 【事实来源】
%   - 上述 crate 源码与 Cargo.toml
%   - os/feature-tree.txt 中指向本 impl 的 feature 链

\paragraph{存储模型}

\code{KlogRingbufInner}维护\code{KLOG\_DESC\_SLOTS}（256）个描述符槽，每槽含\code{KlogRecordMeta}与至多\code{KLOG\_MAX\_RECORD\_BYTES}（1024）字节正文缓冲。槽满时覆盖最旧槽，\code{records\_dropped}递增；若被覆盖序号等于\code{read\_cursor\_seq}，游标前推。\code{KLOG\_TEXT\_RING\_BYTES}（32KiB）供\code{SIZE\_BUFFER}语义；当前实现按槽内缓冲聚合，常量保留供后续纯byte-ring扩展。

\paragraph{追加与迭代}

\code{append}分配单调\code{next\_seq}，超长正文截断并置\code{KlogFlags::TRUNC}。\code{iter\_from(start\_seq)}收集仍存活且\code{seq >= start\_seq}的记录，排序后回调\code{KlogRecordView}。覆盖导致序号缺口时，迭代仅反映仍存活的槽。

\paragraph{读游标语义}

\code{read\_cursor\_seq}指向下一条\code{READ}将返回的序号。\code{peek\_next\_unread}取\code{>= read\_cursor}的最小有效\code{seq}。\code{advance\_read\_cursor}设为\code{after\_seq + 1}，并钳制到\code{oldest\_seq}。\code{clear\_read\_cursor}将游标设为\code{newest + 1}（mark全部已读，物理记录保留供\code{iter\_from}）。

\paragraph{全局锁与中断}

\code{KlogRingbuf}为零大小门面；静态\code{Mutex<Option<KlogRingbufInner>>}。\code{with(f)}进入前\code{disable\_global\_interrupt}，退出恢复，避免与trap并发破坏环一致性。\code{init}清空全局环。

\paragraph{syscall 对接}

聚合\code{klog::syscall::dispatch\_kernel}按action分发：\code{CLOSE}/\code{OPEN}、\code{CONSOLE\_*}为no-op；\code{SIZE\_UNREAD}/\code{SIZE\_BUFFER}查询环；\code{READ}/\code{READ\_CLEAR}/\code{READ\_ALL}经\code{format\_traditional}格式化后由syscall层\code{copy\_to\_user}；WRITE（\code{is\_write\_priority}）写入同一环并置\code{USER}标志。未知action当前\code{panic}。

\paragraph{限制}

\code{CONSOLE\_*}未联动runtime控制台。无\code{/dev/kmsg}设备节点。bring-up不校验uid/\code{CAP\_SYSLOG}。环满覆盖后\code{READ}与\code{iter}对gap的处理以实现为准。
